%!TEX root = ../hutbthesis_main.tex

\chapter{基于飞机辅助感知的自动接近与对接控制方法}

\section{感知任务分析与总体方案设计}

在机场飞机自动牵引作业过程中，牵引车能否安全、准确地完成对接，不仅取决于车辆本体的运动控制能力，还取决于其对目标飞机及作业环境的感知能力。特别是在接近飞机前起落架并执行最终对接时，系统需要同时解决目标身份确认、目标位置识别、姿态估计及近距离安全态势判断等问题。因此，感知系统是实现自动牵引作业闭环控制的关键组成部分\textsuperscript{[18,29-30]}。

结合本课题的研究目标与仿真实现条件，本文引入飞机辅助感知方案，即在飞机前起落架区域设置可识别的特征标识，通过牵引车搭载的视觉与近距离传感器完成目标识别与态势获取。该方案的核心思想是：将原本难以直接观测、难以稳定识别的前起落架对接区域，转化为具有明确标志特征和身份属性的感知目标，使牵引车能够在接近过程中实现“先确认、后接近、再对准”的分阶段感知决策\textsuperscript{[29-30]}。

与传统依赖人工目视判断的作业方式相比，该方案具有以下特点：其一，能够为牵引车提供明确的目标身份信息，避免牵引对象识别错误；其二，能够在接近过程中持续提供目标相对位置与姿态信息，为控制器提供实时反馈；其三，能够将飞机对接区域的局部环境显式化，使碰撞风险、偏航误差和横向偏差能够被量化描述。

因此，飞机辅助感知并不是单纯的“加一个二维码”或“加一个标签”，而是将目标飞机前起落架区域设计为牵引系统可识别、可验证、可跟踪的感知节点，从而为后续自动控制提供稳定的先验信息和观测依据。

基于上述分析，本文构建的感知总体方案由三部分组成：\\
第一部分为目标身份确认模块，用于确认当前牵引任务对应的飞机对象是否与预设目标一致；

第二部分为相对位姿感知模块，用于估计牵引车与飞机前起落架之间的相对距离、横向偏差和姿态偏差；

第三部分为近距态势判断模块，用于在对接过程中识别安全边界、碰撞风险与姿态偏离程度，并向控制层输出约束信息。

从系统结构上看，该方案形成了由“飞机端目标标识---牵引车端感知模块---控制决策节点”构成的信息链路。飞机前起落架不再只是被动的对接对象，而是成为局部识别信息载体；牵引车通过传感器采集这些信息，再将其转换为控制层能够使用的状态量。这样一来，第二章中的车辆运动学模型与本章中的感知输出自然衔接：前者解决牵引车如何运动，后者解决车辆如何获得目标、误差和停车依据。

\section{飞机辅助感知目标与识别信息设计}

为使牵引车能够稳定获取飞机前起落架区域的目标信息，本文将飞机辅助感知对象限定为前起落架附近的局部功能区域，并对该区域引入结构化识别信息。该识别信息并不改变飞机本体的运动学特性，而是作为一种面向自动牵引任务的感知标识，用于提升牵引车对目标区域的可识别性与可确认性。

具体而言，本文在飞机前起落架区域设置特殊识别码，用以表征以下信息：\\
（1）目标飞机身份编号；\\
（2）对接区域中心参考位置；\\
（3）局部方向参考特征；\\
（4）用于视觉识别的几何轮廓特征。

上述信息在仿真环境中可通过纹理标识、平面标记、特征板或特定图案实现，其本质是为牵引车提供一个稳定的局部观测基准\textsuperscript{[29-30]}。与仅依赖起落架自然外形特征相比，这种方式具有识别特征清晰、误识别率低、对光照与视角变化更鲁棒等优点\textsuperscript{[30]}。

在功能上，该识别码作用是身份确认作用。牵引车在进入接近区域后，首先应判断当前视野中的目标是否为任务指定飞机，防止在多机位、多目标环境中出现识别混淆。使牵引车能够从图像观测中计算出目标中心偏移量与方向偏差，从而将视觉结果映射为控制器可直接使用的误差信号。进而实现课题上对“车机互动”的指标要求。

为了表述方便，设飞机辅助感知标识的中心点在飞机局部坐标系中的位置为
\(P_{a}\)，其在牵引车传感器坐标系中的观测位置为
\(P_{s}\)。则通过图像解算或传感器映射后，可得到牵引车与目标之间的相对位置关系。若进一步以对接目标中心作为参考点，则可将识别码的中心点看作前起落架对接区域的等效特征点，后续所有相对偏差计算均围绕该特征点展开。这样做的目的在于将复杂的机械接触区域抽象为一个可观测、可追踪、可控制的数学目标，从而显著降低系统的感知复杂度。

\section{感知系统组成与传感器配置}

结合本课题当前的仿真实现路线，牵引车感知系统采用“视觉识别为主、近距测量为辅”的配置方式。其中，视觉传感器主要负责飞机辅助感知标识的检测与识别，激光雷达或近距传感器主要负责障碍物检测、距离补偿和安全边界判断。该配置既符合现有ROS-Gazebo仿真平台的实现条件，也能够较好覆盖自动接近与对接阶段的主要感知需求。

\subsection{视觉识别模块}

视觉识别模块安装于牵引车前部，主要面向飞机前起落架区域进行图像采集。其核心任务包括：

\begin{quote}
（1）检测比对目标识别码；\\
（2）提取目标区域方向特征；\\
（3）输出目标身份确认结果。
\end{quote}

在实际实现中，该模块可基于深度相机或RGB相机完成\textsuperscript{[29]}。如果采用深度视觉相机，则除可获取二维图像信息外，还可同步获得目标区域的深度信息，有助于提升相对距离估计精度。鉴于本课题当前重点为仿真验证与控制逻辑建立，本文将视觉模块的功能重点放在“目标检测+中心定位+方向识别”三个方面，而不展开复杂的深度神经网络识别框架\textsuperscript{[29-30]}。

\subsection{近距测量模块}

近距测量模块用于补充视觉感知在近距离工况下的局部测量能力。由于牵引车在最终对接阶段距离飞机前起落架较近，此时即使视觉系统能够识别目标，也可能受到视角遮挡、自动对焦的影响。因此，有必要引入激光雷达、超声或等效测距模块，对前方局部区域的距离边界进行测量，以作为视觉估计的校核或补充\textsuperscript{[18]}。

\begin{quote}
该模块的主要作用包括：\\
（1）判断牵引车前端与飞机局部结构之间的最近距离；\\
（2）检测非目标障碍物是否进入安全边界；\\
（3）在最终接近阶段为停车判定提供阈值依据。
\end{quote}

通过视觉与近距测量的组合，牵引车能够在不同距离区间内采用不同感知主导方式：中远距离阶段以视觉为主，近距离阶段以视觉与测距联合判断为主，从而提升系统整体鲁棒性。

\section{目标身份确认模型}

目标身份确认是自动牵引作业的首要感知任务。只有在确认当前目标为预定飞机对象之后，牵引车才应进入自动接近与对接控制阶段。否则，即使轨迹控制精度足够高，也可能因目标对象错误而导致整个作业任务失败。

设任务调度系统下发的目标飞机身份编码为
\(I_{t}\)，牵引车视觉模块从识别码中解算得到的观测身份编码为
\(I_{o}\)。则目标身份确认条件可表示为

\(I_{o} = I_{t}\) （15）

当上式成立时，判定目标身份匹配；当上式不成立时，系统应保持等待状态、重新识别或发出异常警告，而不应直接进入对接控制流程。

上述表达式的物理含义非常明确：它并不是一个几何模型，而是一个任务逻辑判定模型。其作用在于建立从“任务分配”到“目标确认”的第一道闭环。对于机坪多机位、多牵引任务并行的应用场景而言，这一步尤其重要。若无身份确认过程，系统只能依赖“看见一个起落架就去接近”的局部策略，这在真实机场环境中是不可靠的。

进一步地，考虑到识别系统存在误检、漏检或图像丢失的情况，可引入确认置信度
\(C_{I}\)，并设置确认阈值
\(C_{\min}\)。则身份确认判据可扩展为\(C_{I} \geq C_{\min}\)
\textsuperscript{[18,30]}。\(
\)
该式表示：只有当识别结果的置信度高于设定阈值时，系统才认为目标身份识别有效。尽管在当前论文阶段可以不深入展开具体置信度计算方法，但在表述上引入这一概念，有助于增强论文在工程实现层面的严谨性。

\section{相对位姿感知模型}

在完成目标身份确认后，牵引车还需进一步获取自身与飞机前起落架区域之间的相对几何关系。自动接近与对接控制的本质，并不是让牵引车“接近某个图像目标”，而是让其在物理空间中逐渐消除相对位置误差与方向误差。因此，感知系统必须将图像信息转化为可用于控制的相对位姿量\textsuperscript{[18,29-30]}。

为此，定义牵引车参考点为 \(P_{v}\)，飞机辅助感知目标中心为
\(P_{a}\)，则两者之间的相对位置误差向量可表示为

\(\mathbf{e}_{p} = \left\lbrack \begin{array}{r}
e_{x} \\
e_{y}
\end{array} \right\rbrack\ \ \ \) （16）

其中，\(e_{x}\) 表示沿牵引车纵向方向的距离误差，\(e_{y}\)
表示沿牵引车横向方向的偏移误差。若从控制意义出发，\(e_{x}\)
主要决定前进或减速行为，\(e_{y}\) 主要决定转向修正行为。

同时，定义飞机辅助感知目标区域的参考方向角为
\(\theta_{a}\)，牵引车当前航向角为
\(\theta_{v}\)，则相对姿态误差可表示为

\(e_{\theta} = \theta_{a} - \theta_{v}\) （17）

该式表示牵引车当前朝向与目标对接方向之间的夹角偏差。其物理意义在于：即使车辆已经接近目标中心区域，如果
\(e_{\theta}\)仍较大，则机械对接过程仍可能失败。因此，在最终对接阶段，姿态误差控制通常与横向误差控制具有同等重要性。

综合来看，牵引车在自动接近过程中需要实时感知的核心状态量可归纳为

\(\mathbf{z} = \begin{bmatrix}
e_{x} & e_{y} & e_{\theta}
\end{bmatrix}^{T}\) （18）

该状态向量构成本章感知环节向控制环节输出的核心接口。换言之，感知层的任务不是直接控制车辆，而是持续向控制层提供“还差多远、偏了多少、方向还差多少”的量化描述。

\section{近距离态势感知与安全判定}

当牵引车进入飞机前起落架附近的近距离区域后，系统感知重点将从“识别并接近目标”转向“评估局部态势并安全完成对接”。在这一阶段，局部环境中的安全边界、机械接触风险与障碍物侵入风险将显著上升，因此有必要建立近距离态势判断机制。设牵引车前部关键点到飞机局部结构的最小距离为
\(d_{\min}\)，系统设定的安全阈值为
\(d_{s}\)。则基本安全约束可写为\(d_{\min} \geq d_{s}\) 。\(
\)
该式的含义是：在非接触阶段，牵引车与飞机局部结构之间必须始终保持不小于安全阈值的距离，以防止由于感知误差或控制滞后导致提前碰撞\textsuperscript{[5,18]}。

对于最终对接阶段，还可设置停车触发阈值
\(d_{t}\)。当系统判定牵引车与目标对接位置之间的剩余距离满足\(e_{x} \leq d_{t}\)
。\(
\)
且同时满足横向误差与姿态误差均处于允许范围内时，系统可进入停止或低速微调状态。设横向误差允许阈值为
\(e_{y,\max}\)，姿态误差允许阈值为
\(e_{\theta,\max}\)，则对接有效条件可表示为

\(\left\{ \begin{array}{r}
 \mid e_{y} \mid \leq e_{y,max} \\
 \mid e_{\theta} \mid \leq e_{\theta,max}
\end{array} \right.\ \) （19）

这组三重约束共同定义了自动对接的局部成功条件。只有当纵向距离足够接近、横向偏差足够小且方向偏差足够小时，系统才可判定当前已具备机械对接条件。由此可见，近距态势感知并非孤立地“测个距离”，而是要将距离信息、位置误差和姿态误差统一纳入同一判定框架中。

\section{感知流程设计}

结合上述模型，本文将飞机辅助感知系统的工作流程划分为以下步骤：

（1）牵引车进入任务区域后，启动视觉识别模块，对前方目标进行搜索；\\
（2）检测到飞机辅助感知识别码后，解析目标身份信息，并与任务目标编号进行匹配；\\
（3）若身份匹配成功，则提取目标中心点和方向参考特征，计算相对位置误差与姿态误差；\\
（4）结合近距测量模块输出，对当前局部安全状态进行判断；\\
（5）将 \(e_{x}\)、\(e_{y}\)、\(e_{\theta}\)
及安全边界状态发送给控制模块，驱动牵引车完成接近与对准；\\
（6）当对接有效条件满足后，系统进入低速停止或机械连接准备状态。

上述流程体现了“目标确认---误差提取---安全判定---控制反馈”的感知闭环逻辑。与传统单纯依赖图像中心偏差的视觉伺服方法相比，该流程更适合机场牵引这种高安全要求、低速精对位作业任务，因为其将身份确认与近距安全判断一并纳入感知体系，使系统具有更完整的任务语义理解能力。

\section{控制任务描述与总体思路}

在完成第二章牵引车运动学建模以及前文飞机辅助感知方案设计后，系统已经具备了两类基础：其一，双桥转向牵引车的平面运动规律可以用数学模型表示；其二，飞机前起落架区域能够向牵引车提供身份确认结果以及相对位置、相对姿态等感知信息。在此基础上，本节进一步研究牵引车如何利用这些信息完成自动接近与最终对接控制。

从实际作业过程看，牵引车对接飞机前起落架并不是一次性完成的单步动作，而是一个具有明显阶段特征的连续过程。牵引车首先需要确认当前目标是否为任务指定飞机；在确认目标有效后，车辆应由较远位置向飞机前起落架方向稳定接近；当进入近距离区域后，系统控制重点应从“尽快接近目标”转变为“尽快减小横向偏差和姿态偏差”；在最后的微动阶段，车辆需要以极低速度缓慢推进，并在满足安全距离和几何对准条件时稳定停车。因此，自动对接控制的本质不是简单的路径跟踪，而是一个包含任务逻辑、局部轨迹生成、误差反馈控制和终端停车判定的分层闭环控制问题\textsuperscript{[25,27-28]}。

为此，本文采用“模式管理---局部规划---误差反馈---执行器映射”的分层控制方法。其基本思想是：由感知系统实时提供牵引车与飞机辅助感知目标之间的相对误差；由上层逻辑根据当前距离、障碍物情况和识别有效性确定控制模式；由局部规划器生成短时域可行参考轨迹；再由下层控制器根据相对位姿误差计算期望速度和期望转向需求，并进一步映射为前后桥转角与车轮速度命令。该方法既能够保证控制结构清晰，也便于在ROS与Gazebo环境中分模块实现\textsuperscript{[11-16]}。

\section{状态变量、控制输入与误差定义}

\subsection{车辆状态与控制输入}

依据第二章建立的车辆模型，定义牵引车状态向量为

\(\mathbf{x} = \begin{bmatrix}
x & y & \theta & v
\end{bmatrix}^{T}\) （20）

其中，\(x\) 与 \(y\)表示牵引车参考点在全局坐标系中的位置坐标，\(\theta\)
表示车辆航向角，\(v\) 表示车辆纵向线速度。

双桥转向牵引车的控制输入定义为

\(\mathbf{u} = \begin{bmatrix}
v & \delta_{f} & \delta_{r}
\end{bmatrix}^{T}\) （21）

或扩展为

\(\mathbf{u} = \begin{bmatrix}
a & \delta_{f} & \delta_{r}
\end{bmatrix}^{T}\) （22）

其中，\(\delta_{f}\) 与 \(\delta_{r}\)分别为前桥和后桥转角，\(a\)
为车辆纵向加速度。前一种形式适用于直接速度控制实现，后一种形式更适合做速度平滑约束或模型预测控制扩展。

\subsection{感知误差接口}

前文已经将飞机辅助感知结果统一表示为相对误差向量\(\mathbf{z} = \begin{bmatrix}
e_{x} & e_{y} & e_{\theta}
\end{bmatrix}^{T}\)
\(\ \)------\(e_{x}\)表示牵引车与目标区域之间的纵向距离误差；
\(e_{y}\)表示牵引车相对于目标轴线的横向偏差；
\(e_{\theta}\)表示牵引车当前朝向与目标参考方向之间的姿态误差。

在控制意义上，\(e_{x}\) 主要决定车辆应继续前进多少距离，\(e_{y}\)
主要决定车辆应向哪一侧修正，\(e_{\theta}\)
则决定车辆当前朝向是否已经满足对接要求。三者共同构成牵引车自动接近和对准控制的核心反馈量。

\section{对接过程分析与分层控制架构}

\subsection{对接过程分阶段描述}

为了使控制方法与真实作业过程保持一致，本文将牵引车自动对接过程划分为以下五个阶段。

（1）搜索与等待阶段

牵引车进入任务区域后，首先由视觉模块搜索飞机前起落架区域的特殊识别码，并与任务编号进行匹配。若当前识别结果与任务目标不一致，或者识别置信度低于阈值，则系统保持等待状态，不进入自动接近控制。此时输出为零速度、零转角保持命令，以防止误接近。

（2）自动接近阶段

当目标身份确认有效后，系统开始进入自动接近阶段。该阶段的主要目标是让牵引车从较远初始位置稳定接近飞机前起落架前方的预对接区域。此时控制重点不在于最终精确对位，而在于快速减小纵向距离，同时保持可接受的横向误差和姿态误差。

（3）姿态对准阶段

当牵引车与目标之间的纵向距离降低到一定阈值后，系统切换到姿态对准阶段。该阶段中，车辆不再追求较高接近速度，而是优先减小横向误差
\(e_{y}\)和姿态误差
\(e_{\theta}\)，使牵引车中心线逐步与飞机前起落架对接方向重合。

（4）微动对接阶段

当姿态误差和横向误差已经较小，且车辆进入最终接近区域后，系统进入微动对接阶段。此时车辆以低速蠕动方式向前推进，控制重点从“大范围修正”转为“小范围精调”，避免出现冲撞、擦碰或越位现象。

（5）停车保持阶段

当牵引车与目标区域之间的纵向、横向和姿态误差同时满足终端条件，且安全距离约束有效时，系统进入停车保持阶段。此时发布零速度命令，并将前后桥转角回中，保持车辆稳定等待后续机械连接或任务切换。

\subsection{分层控制架构}

根据上述阶段分析，本文建立的分层控制架构如图3-1所示。

这种架构的优点在于：可以把任务逻辑与连续控制分离，从而降低算法复杂度，并且更加符合ROS节点化实现习惯。

\begin{figure}[htbp]
    \centering
    \includegraphics[width=0.82\linewidth]{image6.png}
    \caption{分层控制架构}
\end{figure}

\section{误差动力学模型建立}

为了构造适用于对接过程的控制律，需要首先建立牵引车相对于参考轨迹或参考目标点的误差动力学模型。

设局部参考状态为\(\mathbf{q}_{r} = \begin{bmatrix}
x_{r} & y_{r} & \theta_{r}
\end{bmatrix}^{T}\)。

其中，\(\left( x_{r},y_{r} \right)\) 为局部参考点位置，\(\theta_{r}\)
为参考航向。进一步定义参考速度和参考角速度分别为 \(v_{r}\)与
\(\omega_{r}\)。

将全局位置误差变换到车体坐标系，可以得到

\(\left\{ \begin{array}{r}
\left\lbrack \begin{array}{r}
e_{x} \\
e_{y}
\end{array} \right\rbrack = \begin{bmatrix}
\cos\theta & \sin\theta \\
 - sin\theta & \cos\theta
\end{bmatrix}\left\lbrack \begin{array}{r}
x_{r} - x \\
y_{r} - y
\end{array} \right\rbrack \\
e_{\theta} = \theta_{r} - \theta
\end{array} \right.\ \) (23)

该式的作用是把全局误差转化为车体系误差，使纵向和横向误差具有明确物理意义。其中，\(e_{x}\)
表示车辆前方还剩多少距离，\(e_{y}\) 表示车辆当前偏离对接轴线的程度。

对上述误差求导，并代入车辆运动学方程，可得误差动力学：

\(\left\{ \begin{array}{r}
{\dot{e}}_{x} = v_{r}cose_{\theta} - v + \omega e_{y} \\
{\dot{e}}_{y} = v_{r}sine_{\theta} - \omega e_{x} \\
{\dot{e}}_{\theta} = \omega_{r} - \omega
\end{array} \right.\ \) (24)

上述误差方程表明：

\begin{enumerate}
\def\labelenumi{\arabic{enumi}.}
\item
  纵向误差 \(e_{x}\)的变化不仅与前进速度
  \(v\)有关，还受姿态误差和横向运动耦合影响；
\item
  横向误差 \(e_{y}\)的变化主要由角速度控制；
\item
  姿态误差 \(e_{\theta}\)的变化取决于参考角速度与实际角速度的差异。
\end{enumerate}

这说明自动对接控制本质上是一个非线性耦合误差调节问题，不能简单把纵向和横向完全独立看待\textsuperscript{[19,27-28]}。

\section{局部参考轨迹生成方法}

\subsection{预对接点与对接点构造}

为了避免车辆直接朝最终机械连接位置冲击，本文采用两级目标点设计，即“预对接点”和“对接点”。

设飞机辅助感知系统给出的目标中心点为

\(\mathbf{p}_{a} = \left\lbrack \begin{array}{r}
x_{a} \\
y_{a}
\end{array} \right\rbrack\) (25)

目标方向为 \(\theta_{a}\)。则预对接点定义为

\(\mathbf{p}_{p} = \mathbf{p}_{a} - d_{pre}\left\lbrack \begin{array}{r}
\cos\theta_{a} \\
\sin\theta_{a}
\end{array} \right\rbrack\) （26）

\(d_{pre}\)
为预留接近距离。该点位于最终目标点前方一定距离处，主要用于完成姿态调整。最终对接点定义为

\(\mathbf{p}_{d} = \mathbf{p}_{a} - d_{hook}\left\lbrack \begin{array}{r}
\cos\theta_{a} \\
\sin\theta_{a}
\end{array} \right\rbrack\) （27）

\(d_{hook}\)
为机械连接前的停车距离。该点用于微动对接阶段的最终停靠。如果牵引车从较远位置直接朝最终对接点收敛，往往会导致中后段横向修正不足或姿态调节过猛；引入预对接点后，车辆可以先对准姿态，再进行最后的缓慢接近\textsuperscript{[20-21,26-28]}。

\section{参考速度调度策略}

在机场牵引对接作业中，恒定速度控制并不合适。当路径弯曲程度较大时，高速运行会导致横向误差难以收敛；
而当障碍物距离较小时，必须主动减速保证停车裕度；
当进入微动阶段时，也需要明显降低速度以提高最终对位精度。
因此，参照相关文献资料，认为可以采用如下速度调度公式：

\(v_{r} = min\left( v_{stage},\sqrt{\frac{a_{y,max}}{\mid \kappa_{r} \mid + \varepsilon_{\kappa}}},k_{d}\lbrack d_{\min} - d_{s}\rbrack_{+} \right)\)
(28）

其中：\(v_{stage}\)为当前控制模式允许的最大速度；
\(a_{y,\max}\)为允许的最大横向加速度； \(\kappa_{r}\)为参考轨迹曲率；
\(d_{\min}\)为前方最近障碍距离； \(d_{s}\)为安全距离阈值；
\(k_{d}\)为距离---速度缩放系数\textsuperscript{[25,27-28]}。

\section{误差反馈控制律设计}

\subsection{纵向速度控制律}

纵向控制律设计为

\(v_{d} = {sat}_{\left\lbrack 0,v_{\max} \right\rbrack}\left( v_{r}cose_{\theta} + k_{x}e_{x} \right)\)
(29）

其中，\(k_{x}\) 为纵向误差增益，\(sat\) 表示饱和函数。

该控制律由两部分组成：\(v_{r}\cos e_{\theta}\)：表示当车辆姿态误差较大时，自动减小前进速度，防止车辆在未对准状态下强行推进；
\(k_{x}e_{x}\)：表示根据纵向误差大小补偿推进速度，使车辆逐步靠近目标。

这种设计符合车辆对接实际过程：距离远时可较快前进；姿态偏差大时必须主动减速\textsuperscript{[19]}。

\subsection{角速度控制律}

角速度控制律设计为

\(\omega_{d} = {sat}_{\left\lbrack - \omega_{\max},\omega_{\max} \right\rbrack}\left( \omega_{r} + v_{r}(k_{y}e_{y} + k_{\theta}sine_{\theta}) \right)\)
（30）

其中，\(k_{y}\) 和
\(k_{\theta}\)分别为横向误差增益和姿态误差增益。该控制律的物理意义如下\textsuperscript{[17,26-28]}：

\(\omega_{r}\)：参考轨迹自身的前馈角速度，用于保证车辆沿局部参考曲线运行；

\(v_{r}k_{y}e_{y}\)：根据横向偏差调节转向，使车辆向目标轴线靠拢；

\(v_{r}k_{\theta}\sin e_{\theta}\)：根据姿态误差修正车辆方向。

之所以采用 \(\sin e_{\theta}\)而不是直接用
\(e_{\theta}\)，是因为前者在大角度误差情况下更平滑，且在小角度时近似等于
\(e_{\theta}\)，便于后续稳定性分析。

\subsection{对接过程中的控制行为分析}

结合车辆对接全过程，上述控制律在不同阶段表现如下：

\begin{quote}
（1）自动接近阶段：\\
此时 \(e_{x}\)较大，速度项
\(k_{x}e_{x}\)占主导，车辆以较快速度向预对接点推进。
\end{quote}

（2）姿态对准阶段\\
随着 \(e_{x}\)逐步减小，横向误差 \(e_{y}\)和姿态误差
\(e_{\theta}\)开始主导控制器输出，角速度命令增大，车辆重点完成方向修正。

（3）微动对接阶段\\
此时
\(v_{r}\)已被速度调度压低，车辆角速度输出较平缓，控制器主要进行细小偏差修正，避免出现末端振荡。

因此，这组控制律不仅在数学上闭合，而且在物理意义上能够与整个对接过程自然对应。

\section{四轮转向映射与约束处理}

\subsection{曲率计算}

由期望线速度和角速度得到期望曲率

\(\kappa_{d} = \frac{\omega_{d}}{max(v_{d},\varepsilon_{v})}\) （31）

其中, \(\varepsilon_{v}\)为防止低速时分母趋近于零的安全常数。

\subsection{双桥反向转向模式}

在自动接近阶段，如果需要较大机动能力，则采用双桥反向转向模式\textsuperscript{[19,28]}：

\(\left\{ \begin{array}{r}
\delta_{f}^{\star} = sat\left( arctan\frac{L\kappa_{d}}{2}, - \delta_{f,max},\delta_{f,max} \right) \\
\delta_{r}^{\star} = - \delta_{f}^{\star}
\end{array} \right.\ \) (32)

该模式可以显著减小等效转弯半径，适合从较远位置快速接近飞机前起落架区域。

\subsection{前桥单转向模式}

在姿态对准和微动阶段，为了提高控制平稳性，采用前桥单转向模式：

\(\left\{ \begin{array}{r}
\delta_{f}^{\star} = sat\left( arctan(L\kappa_{d}), - \delta_{f,max},\delta_{f,max} \right) \\
\delta_{r}^{\star} = 0
\end{array} \right.\ \) (33)

该模式虽然机动性不如双桥反向转向，但执行器协调更简单，转向响应更平顺，适合最终精调。

\subsection{速度与转角限幅}

为了保证控制命令可执行，还需对速度与转角进行约束处理：

\(\left\{ \begin{array}{r}
0 \leq v_{d} \leq v_{\max} \\
 \mid \delta_{f}^{\star} \mid \leq \delta_{f,max} \\
 \mid \delta_{r}^{\star} \mid \leq \delta_{r,max}
\end{array} \right.\ \) (34)

同时对速度变化率和转角变化率限制为

\(\left\{ \begin{array}{r}
 \mid v_{d}(k) - v_{d}(k - 1) \mid \leq a_{\max}T_{s} \\
 \mid \delta_{i}(k) - \delta_{i}(k - 1) \mid \leq {\dot{\delta}}_{\max}T_{s},i \in \{ f,r\}
\end{array} \right.\ \) (35)

这样可以避免由于离散控制导致的命令突变，从而减小Gazebo中常见的抖振与过冲现象。

\section{终端停车判定与模式切换}

\subsection{终端停车集合}

在最终微动阶段，当牵引车同时满足以下条件时，可判定已具备停车条件：

\(\left\{ \begin{array}{r}
\begin{array}{r}
 \mid e_{x} \mid \leq d_{t} \\
 \mid e_{y} \mid \leq e_{y,max}
\end{array} \\
 \mid e_{\theta} \mid \leq e_{\theta,max} \\
d_{\min} \geq d_{s}
\end{array} \right.\ \) (36)

将其定义为终端停车集合

\(\Omega_{t} = \left\{ \mid e_{x} \mid \leq d_{t},\text{ } \mid e_{y} \mid \leq e_{y,max},\text{ } \mid e_{\theta} \mid \leq e_{\theta,max},\text{ }d_{\min} \geq d_{s} \right\}\)
(37)

与单纯依靠纵向距离判定停车不同，这种集合式判定可以避免“距离够近但方向没对准”的错误停车问题。

\subsection{模式切换规则}

系统模式切换规则如下：

搜索-自动接近：识别码身份确认成功且置信度满足阈值；

自动接近-姿态对准：\(e_{x}\) 进入近距阈值；

姿态对准-微动对接：\(e_{x}\) 足够小且 \(e_{\theta}\)已较小；

任意模式-中止：目标丢失或最近距离低于中止阈值；

微动对接-停车保持：进入终端集合 \(\Omega_{t}\)。

这意味着控制系统并不是依赖一条连续控制律从头走到尾，而是通过“阶段切换 +
连续控制”的混合方式实现更符合工程需求的动作组织。

\section{控制算法主流程}

图3-2给出自动接近与对接控制算法主流程。

\begin{figure}[htbp]
    \centering
    \includegraphics[width=0.82\linewidth]{image7.png}
    \caption{全过程算法流程图}
\end{figure}

\section{本章小结}

本章将原有飞机辅助感知与自动接近控制内容整合为同一技术链路。首先，围绕飞机前起落架对接区域设计了目标标识、视觉识别和近距测量方案，并给出了目标身份确认、相对位姿误差和安全边界判定方法；其次，在第二章双桥转向运动学模型基础上，建立了车体系误差动力学模型，设计了预对接点、对接点、速度调度、误差反馈控制和四轮转向映射方法；最后，通过终端停车集合与模式切换规则，将目标确认、接近、姿态对准、微动对接和停车保持组织为完整闭环。

与单纯的路径跟踪控制相比，本章方法更加贴合机场低速高精度对接工况，能够把目标身份确认、局部态势感知、轨迹可行性、误差收敛性和停车安全性放在同一控制框架中，为后续仿真实验与系统实现提供算法基础。
